\begin{table*}[tp]
  \centering
  \caption{Adjacent Pair Accuracy by fragment count. Difficulty scales with how many pieces a protein was digested into: the number of possible orderings grows factorially while the evidence available at each junction does not. Lift is the mean paired per-sample difference between the LLM-Guided arm and the Random Order floor within the bin, not a difference of two independent means. Setup: \textnormal{\textnormal{\textit{S. cerevisiae}}}, 5 digestion replicas.}
  \label{tab:stratification_yeast_r5}
  \begin{tabular}{lrrrrrr}
    \toprule
    Fragments & n & Random Order & Fixed Settings & Random Search (no LLM) & LLM-Guided & Lift \\
    \midrule
    2-4 & 4 & 0.500 & 0.667 & 1.000 & 1.000 & +0.500 \\
    5-9 & 18 & 0.076 & 0.446 & 0.576 & 0.568 & +0.492 \\
    10-19 & 21 & 0.092 & 0.311 & 0.409 & 0.431 & +0.339 \\
    20-49 & 37 & 0.037 & 0.185 & 0.295 & 0.310 & +0.274 \\
    50+ & 20 & 0.023 & 0.147 & 0.227 & 0.248 & +0.225 \\
    \bottomrule
  \end{tabular}
  \par\smallskip\footnotesize{n counts proteins in the bin; a bin's mean is taken over those whose ordering metrics are defined (true fragment order recovered), so it can rest on fewer than n.}
\end{table*}
